\section{Introduction}
\label{sec:intro}

Floods are among the most frequent and damaging natural hazards, and timely
inundation maps are needed most while the storm is still overhead. Optical
imagery is often unusable at exactly that moment because of cloud cover.
Synthetic aperture radar (SAR) does not have this limitation. The Sentinel-1
mission delivers C-band SAR acquisitions day and night and through cloud
\cite{torres2012}. Open water often appears dark because smooth water surfaces
scatter the signal away from the sensor \cite{twele2016sentinel1}. This has
made Sentinel-1 a workhorse of operational flood mapping.

Deep segmentation networks trained on the Sen1Floods11 benchmark
\cite{bonafilia2020} are a strong data-driven alternative to threshold-based
processing chains. The standard evaluation, however, is a chip-level split
in which training and test chips can originate from the same flood events.
Scores obtained this way describe how well a model interpolates within
events it has already seen. What a practitioner needs to know is whether
the model generalizes to an event that contributed no training labels at
all. This matters for regions such as Azerbaijan that are absent from
Sen1Floods11.

We therefore use a leave-one-event-out protocol, retaining the historical
LOCO (leave-one-country-out) identifier. Each Sen1Floods11 hand-labeled
group corresponds to a flood event, so an excluded group represents a
particular event and region. A raw held-out IoU is hard to interpret on its
own, because some events are simply harder than others. We compare the
available event-excluded checkpoints with a fixed generalist on exactly
the same official test chips. This pairs the evaluation inputs while
leaving differences between the training and validation designs
(Sec.~\ref{sec:paired}).

Segmentation alone is also insufficient for a usable inundation map.
Wind-roughened water can look like land, recurrent water bodies look like
flood water, and SAR terrain distortions such as layover and shadow can
produce false detections. We address the last two with a post-processing
fusion that combines the JRC Global Surface Water layer \cite{pekel2016high}
with a DEM-derived terrain and coverage mask, separately from the
benchmark (Sec.~\ref{sec:fusion}).

\paragraph{Contributions.}
\begin{itemize}
\item We reconstruct a paired, official-test comparison on seven
Sen1Floods11 events and nine recoverable event--seed combinations
(Secs.~\ref{sec:dataset} and \ref{sec:loco}).
\item We report the seed-specific difference and show how it changes
when individual test chips are omitted (Sec.~\ref{sec:sensitivity}).
\item We document the VV/VH/ratio input, weak-to-hand training, and an
11-variant loss sweep under class imbalance
(Secs.~\ref{sec:method} and \ref{sec:ablations}).
\item We describe the prototype's recurrent-water and terrain/coverage
filters and their effect on two unlabeled Azerbaijani scenes
(Sec.~\ref{sec:deployment}).
\end{itemize}

\section{Related Work}
\label{sec:related}

\paragraph{SAR flood mapping.}
Classical Sentinel-1 flood mapping relies on automated thresholding and
change detection chains \cite{twele2016sentinel1}. Learning-based methods
became practical with georeferenced training data. Sen1Floods11
\cite{bonafilia2020} provides hand-labeled and weakly labeled Sentinel-1
chips from flood events on several continents and is the benchmark we
build on.

\paragraph{Ancillary layers and terrain effects.}
The JRC Global Surface Water dataset \cite{pekel2016high} records how often
each pixel has been observed as water over several decades and provides a
prior for recurrent water. On the SAR side, steep terrain produces layover
and shadow, whose extent depends on the viewing geometry and the local
slope. The same geometry underlies radiometric terrain correction
\cite{small2011}. Elevation data such as the Shuttle Radar Topography
Mission \cite{farr2007} make these effects computable per pixel.

\paragraph{Evaluating generalization.}
Random splits of spatially autocorrelated data can overestimate
performance on new locations, motivating blocked or leave-one-group-out
validation \cite{roberts2017}. T{\"u}rkmenli, Kayabol and Aptoula
\cite{turkmenli2025} treat flood events as source domains for
domain-generalization training, whereas we compare an event-excluded
checkpoint with a fixed generalist on that event's same official test
chips.

\paragraph{Segmentation under class imbalance.}
Water pixels are a small minority. Overlap-based objectives such as Dice
\cite{milletari2016}, Tversky \cite{salehi2017} and Focal-Tversky
\cite{abraham2019}, the Lov\'asz surrogate for IoU \cite{berman2018},
and focal loss \cite{lin2017focal} are common answers. We compare them
with weighted cross-entropy in a validation sweep (Sec.~\ref{sec:ablations}).
